% ============================================================
\chapter{まとめ}\label{ch:summary}

この資料は、Black-Litterman モデルのビュー $Q$、その不確かさ $\Omega$、均衡の不確かさ $\tau$ の決め方を、論文を読んで整理し、本プロジェクトのデータと論文の数値例で確かめた。

\section{分かったこと}

\begin{description}[leftmargin=2em,style=nextline]
\item[1. BL の結果は採用率で決まる] 均衡からビューへ何割動くかは、$\Omega$ と $\tau\Sigma$ の比で決まる。He--Litterman の既定値では 50\%、予想の当たり具合に合わせると IC の 2 乗（多くは 1\% 未満）で、この違いが BL の結果の違いの大部分を生む。$\Omega$ を $\tau$ に比例させれば $\tau$ は消え、$\Omega$ を外から決めれば $\tau$ が採用率の係数になる。
\item[2. 確信度の方法は閉じた式になる] Idzorek の 8 資産の例は論文の表どおりに再現でき、確信度の方法は $\omega_k=\tau p_k\Sigma p_k\T(1-C_k)/C_k$ と同じだった。確信度は 1 つのビューの採用率で、ビューが複数あると相関によって確信度どおりには動かない。
\item[3. 当たり具合から $\Omega$ を決めると、Grinold の式になる] 市場版の BL で $\omega=\sigma^2(1/\IC^2-1)$ と置くと、事後の期待リターンは $\IC\cdot\sigma\cdot z$ に一致する。BL は、昔からの「予想を当たり具合で縮める」実務と同じものである。
\item[4. 対角の $\Omega$ は確信を積み上げる] 同じ情報源から多くのビューを作ると、外れには共通の部分がある。$\Omega$ を対角にすると、それを無視して、ビューの数だけ確信が積み上がる。本プロジェクトのデータで、1 銘柄なら 5〜7\% の採用率が、対角の $\Omega$ では中央値 \V{bt.calib.adopt}\% になった。
\item[5. 道具が出す不確かさは、当たり具合と合わない] LLM の回答のばらつきは、実際の外れの約 8 分の 1 だった。調査メモの「ばらつきが不確かさを客観的に表す」という主張は、本プロジェクトのデータでは成り立たない。
\item[6. ファクターごとの材料は、使い方を選ぶ] 政治の指数は、短い期間の不確かさでは市場のボラティリティに及ばず、シナリオや長い期間の見込みに向く。マクロのモデルの平均の式は危機で極端になりうる。ESG を共分散に入れる案は、入れ方によって効かないか、意図と逆に働く。モメンタムは、自身のボラティリティで確信度を下げると急落を和らげられる。
\item[7. 拡張は、平均以外のビューと裾のリスクのため] エントロピー・プーリングは、平均のビューなら BL とほぼ同じ答えを出し、順位・ボラティリティ・下落の確率のビューも入れられる。CVaR は、楕円分布に近い資産では分散と同じ答えになる。
\item[8. 最適化ではずれの大きさを別に決める] 正しく縮めた期待リターンでも、平均分散最適化は大きなトラッキングエラーを取る。$\Omega$ は当たり具合で、TE は制約で決める。事前の TE は実績より小さく出やすい。
\end{description}

\section{本プロジェクトへの示唆}

本プロジェクトの BL は、いまは LLM の回答の分散を $\Omega$ にしており、採用率が非常に高く、比率は少数の銘柄に集中する。この資料の結果から、次の順に直すのがよいと考える（\chapref{ch:project}）。

\begin{enumerate}
\item 予想と結果から、判断日ごとに較正の傾き $b_t$ を計算する。
\item $\Omega=\tau\Sigma(1/b_t-1)$ とし、採用率を $b_t$ にそろえる。
\item 予想から判断日ごとの平均を引き、相対ビューにする。
\item 事前の TE に上限を置き、実績の TE を見て調整する。
\item 採用率、実質銘柄数、TE、IC を毎回記録し、成績だけで方法を選ばない。
\end{enumerate}

これらは、LLM の予想が「当たるかどうか」を判断する前に必要な、測り方の整備である。予想の価値（IC が本当に 0 より大きいか）は、記録をためて確かめるしかない。BL は、価値の分からない予想を、分からない程度に応じて控えめに使う仕組みとして働かせるのがよい。

\section{残った問い}

\begin{itemize}
\item LLM の予想の当たり具合は、局面（上昇相場・下落相場、ボラティリティの高さ）でどう変わるか。変わるなら、$\Omega$ をどう動かすべきか。
\item 複数の LLM の予想の外れは、どれくらい相関しているか。融合で本当に改善するか。
\item セクターやファクターのような、数の少ないビューにしたとき、LLM の予想の当たり具合は上がるか。
\item 文章の情報（月次レポート、決算説明会）を加えたとき、IC は上がるか。
\end{itemize}

どの問いも、この資料で示した測り方（IC、較正の傾き、外れと不確かさの比、採用率、TE）で答えられる。
